
\section{10. Limitations}

\subsection{10.1 Task Coverage}

Only SST-2 evaluated ($L = 128$). Primal Attention's advantage grows at 
longer $L$; not tested at $L \ge 2048$.

\subsection{10.2 Model Coverage}

Tested on models with $H = 12$. Not tested on:
\begin{itemize}
\item Very wide models ($H = 32+$)
\end{itemize}
\begin{itemize}
\item Mixture-of-Experts
\end{itemize}
\begin{itemize}
\item Vision transformers
\end{itemize}

\subsection{10.3 Approximate Causality}

Causal masking is applied to the approximate attention matrix. This is 
not equivalent to true causal attention. We do not measure the impact on 
autoregressive generation.

\subsection{10.4 No True Speedup at Short $L$}

At $L = 128$, Primal Attention is {{D\_TIME\_RATIO}}$\times$ slower. The overhead of 
computing three factors ($F_1, F_2, F_3$) plus pseudoinverse exceeds the 
savings from reduced attention matrix.

\subsection{10.5 Segment Pooling Bias}

Segment pooling is content-agnostic. Attention sinks and heavy hitters are 
not explicitly preserved.

\subsection{10.6 No Parameter Reduction}

Primal Attention does \textbf{not} reduce parameter count. For storage savings, 
use compression methods (A3, AASVD).

\subsection{10.7 Single Seed}

All results from seed {{SEED}}.

\subsection{10.8 No Comparison with Performer / Linformer}

We do not directly compare with other linear attention methods. This limits 
external validity.

\subsection{10.9 Failure Modes}

{{UNRESOLVED\_FAILURES}}

\subsection{10.10 Code Quality}

Research-grade implementation. Some numerical stability issues may remain 
at extreme $L$ or $m$.

\subsection{10.11 Ethical Considerations}

Dual-use. See Appendix E.